% =====================================================================
\section{Related Work}\label{sec:related}
% =====================================================================

\subsection{Arrival-Time Prediction for Collection Trucks and Buses}
Chen \emph{et al.}~\cite{chen2019travel} predicted waste-truck travel
times in Hsinchu City, Taiwan, by clustering historical inter-stop trips
and fitting a linear regression per cluster. The clustered model
outperformed unclustered regression and a neural network, and the authors
favored it partly for its low computation cost. Their app lets a user
pick a nearby point and view its predicted arrival time, but the paper
does not consider whether the user can walk there in time. For buses, Zhang
\emph{et al.}~\cite{zhang2024bus} tuned an LSTM with an improved whale
optimization algorithm on one day of data from one route, which limits
generalization. Rawat \emph{et al.}~\cite{rawat2026bus} trained a
route-embedded LSTM on 20 New York City bus routes and reported a mean
absolute error (MAE) of about 2.02 minutes, a 28.5\% improvement over a
historical-average baseline. Buses, however, differ from garbage trucks
in dwell behavior and stop purpose. Taichung's own system assumes that a
delay observed at one point propagates unchanged to all later points and
updates its estimates every minute~\cite{tcepb2022report}. Its contract
requires an error within $\pm5$ minutes, but no measured error is
published.

\subsection{Walking Time and Pedestrian Routing}
Fink \emph{et al.}~\cite{fink2024travel} built a travel-time matrix for
Helsinki with separate walking speeds for average and slower walkers,
showing that one speed can overstate how far slower people can go. Al
Shammas \emph{et al.}~\cite{alshammas2023gis} built a GIS pedestrian
network of sidewalks and crossings for comfort-aware routing in Madrid.
Neither study combines walking time with a vehicle ETA.

\subsection{LLM Agents and Tool Calling}
Toolformer~\cite{schick2023toolformer} showed that language models can
learn to call external APIs, and ReAct~\cite{yao2023react} interleaves
reasoning with tool actions. LLMAP~\cite{yuan2025llmap} uses an LLM only
to parse user needs into structured constraints and leaves route search
to a graph algorithm, which outperformed LLM-as-agent planning. The
authors also found that LLMs estimate numerical preference weights
poorly. These findings motivate our rule that the LLM never produces ETA,
distance, or walking-time values itself.

\subsection{Waste Collection Systems and Open Data Quality}
Most smart-collection studies serve operators: GIS route optimization in
Dubai~\cite{hassan2025dubai}, a meta-analysis of IoT routing that found a
21.51\% average distance reduction (only 12.37\% with real operating
data)~\cite{maciel2025iot}, and waste-volume patterns in
Bangkok~\cite{hamontree2026bangkok}. On the resident side, a Lima system
sent proximity alerts when the truck approached, and residents who knew
the collection time rose from 30\% to 95\%~\cite{duran2026lima}, but it
did not predict arrival or consider walking time. A survey of 215 data
stewards on Brazil's open data portal found that only 54\% updated data
at the declared frequency~\cite{tejedo2025brazil}, which shows that open
data timeliness cannot be assumed.

\subsection{Research Gap}
Table~\ref{tab:comparison} compares GarbageGo with the closest systems.
Within the work we reviewed, no resident-facing system combines truck ETA,
walking time, and data uncertainty to decide whether a user can make it
now, or offers an alternative when the user cannot. GarbageGo targets
this gap under the constraint of low-frequency, incomplete open data.

% =====================================================================
% Table: Capability comparison with existing resident-facing services
% Included from sections/02_related_work.tex via % =====================================================================
% Table: Capability comparison with existing resident-facing services
% Included from sections/02_related_work.tex via % =====================================================================
% Table: Capability comparison with existing resident-facing services
% Included from sections/02_related_work.tex via \input{tables/comparison}
% "--" = not mentioned in the public document or paper (not proof of absence)
% =====================================================================
\begin{table}[htbp]
\caption{Capability Comparison With Existing Services}
\label{tab:comparison}
\centering
\setlength{\tabcolsep}{3pt}
\footnotesize
\begin{tabular}{@{}lcccc@{}}
\toprule
\textbf{Capability} & \textbf{Taichung~\cite{tcepb2021ethics}} &
\textbf{Hsinchu~\cite{chen2019travel}} & \textbf{Lima~\cite{duran2026lima}} &
\textbf{Ours} \\
\midrule
Real-time truck position   & \checkmark & \checkmark & \checkmark & \checkmark \\
Nearby point search        & \checkmark & \checkmark & --         & \checkmark \\
Arrival-time prediction    & $\circ$    & \checkmark & --         & \checkmark \\
Walking time               & --         & --         & --         & \checkmark \\
Can-I-make-it judgment     & --         & --         & --         & \checkmark \\
Alternative suggestion     & --         & --         & --         & \checkmark \\
Data-freshness warning     & --         & --         & --         & \checkmark \\
\bottomrule
\multicolumn{5}{@{}p{0.95\columnwidth}@{}}{\scriptsize
\checkmark: provided (Ours: planned); $\circ$: mentioned, method not
disclosed; --: not mentioned in the public document.}
\end{tabular}
\end{table}

% "--" = not mentioned in the public document or paper (not proof of absence)
% =====================================================================
\begin{table}[htbp]
\caption{Capability Comparison With Existing Services}
\label{tab:comparison}
\centering
\setlength{\tabcolsep}{3pt}
\footnotesize
\begin{tabular}{@{}lcccc@{}}
\toprule
\textbf{Capability} & \textbf{Taichung~\cite{tcepb2021ethics}} &
\textbf{Hsinchu~\cite{chen2019travel}} & \textbf{Lima~\cite{duran2026lima}} &
\textbf{Ours} \\
\midrule
Real-time truck position   & \checkmark & \checkmark & \checkmark & \checkmark \\
Nearby point search        & \checkmark & \checkmark & --         & \checkmark \\
Arrival-time prediction    & $\circ$    & \checkmark & --         & \checkmark \\
Walking time               & --         & --         & --         & \checkmark \\
Can-I-make-it judgment     & --         & --         & --         & \checkmark \\
Alternative suggestion     & --         & --         & --         & \checkmark \\
Data-freshness warning     & --         & --         & --         & \checkmark \\
\bottomrule
\multicolumn{5}{@{}p{0.95\columnwidth}@{}}{\scriptsize
\checkmark: provided (Ours: planned); $\circ$: mentioned, method not
disclosed; --: not mentioned in the public document.}
\end{tabular}
\end{table}

% "--" = not mentioned in the public document or paper (not proof of absence)
% =====================================================================
\begin{table}[htbp]
\caption{Capability Comparison With Existing Services}
\label{tab:comparison}
\centering
\setlength{\tabcolsep}{3pt}
\footnotesize
\begin{tabular}{@{}lcccc@{}}
\toprule
\textbf{Capability} & \textbf{Taichung~\cite{tcepb2021ethics}} &
\textbf{Hsinchu~\cite{chen2019travel}} & \textbf{Lima~\cite{duran2026lima}} &
\textbf{Ours} \\
\midrule
Real-time truck position   & \checkmark & \checkmark & \checkmark & \checkmark \\
Nearby point search        & \checkmark & \checkmark & --         & \checkmark \\
Arrival-time prediction    & $\circ$    & \checkmark & --         & \checkmark \\
Walking time               & --         & --         & --         & \checkmark \\
Can-I-make-it judgment     & --         & --         & --         & \checkmark \\
Alternative suggestion     & --         & --         & --         & \checkmark \\
Data-freshness warning     & --         & --         & --         & \checkmark \\
\bottomrule
\multicolumn{5}{@{}p{0.95\columnwidth}@{}}{\scriptsize
\checkmark: provided (Ours: planned); $\circ$: mentioned, method not
disclosed; --: not mentioned in the public document.}
\end{tabular}
\end{table}

